\begin{exercisechunk}
\item \textbf{A small teaching set for a maximum-margin learner.}\suppmark
\label[exercise]{ex:l22-max-margin-teaching}
\exerciseprofile{3}{3}{3}{\exproof\quad\exconstruction\quad\exinterpretation}{%
Use active constraints to teach a specified target to a specified learner.}
Let $\mathcal X\subset\mathbb R^p$ be finite, and let
$f:\mathcal X\to\{-1,1\}$ be affine-separable with both labels present.
The full labeled dataset is
$D_0=\{(x,f(x)):x\in\mathcal X\}$. For a labeled subset $D$, define
$D_+=\{x:(x,1)\in D\}$ and $D_-=\{x:(x,-1)\in D\}$. The learner chooses
\[
 w_D=\arg\min_w\frac12\lVert w\rVert_2^2
 \quad\text{subject to}\quad
 \langle w,x^+-x^-\rangle\ge2
 \quad\text{for every }x^+\in D_+,\ x^-\in D_-,
\]
and then sets
\[
 b_D=-\frac12\left(
 \min_{x^+\in D_+}\langle w_D,x^+\rangle
 +\max_{x^-\in D_-}\langle w_D,x^-\rangle\right).
\]
The returned rule is $x\mapsto\operatorname{sign}(\langle w_D,x\rangle+b_D)$.
You may use the conic Carath\'eodory theorem: every vector in a cone generated
in $\mathbb R^p$ can be expressed as a nonnegative combination of at most
$p$ generators. Using the weaker bound $p+1$ is also acceptable.
\begin{enumerate}[(a)]
\item Let $w_0=w_{D_0}$. Use the first-order optimality condition to show
that $w_0$ is a nonnegative combination of at most $p+1$ active difference
vectors $x_j^+-x_j^-$.
\item Let $D$ contain the endpoints of the selected active differences.
Prove $|D|\le2p+2$ and $w_D=w_0$.
\item Prove that the selected positive endpoints attain
$\min_{x^+\in(D_0)_+}\langle w_0,x^+\rangle$ and that the selected negative
endpoints attain
$\max_{x^-\in(D_0)_-}\langle w_0,x^-\rangle$. Deduce $b_D=b_{D_0}$ and
prove that the learner returns $f$ on all of $\mathcal X$.
\item Explain why the construction does not provide one universal dataset
that teaches every affine rule or every learning algorithm.
\end{enumerate}
\begin{hints}
The normal cone of the feasible polyhedron at $w_0$ is generated by the
negatives of its active constraint normals. For (c), pair a hypothetical
smaller positive projection with a selected negative endpoint.
\end{hints}
\end{exercisechunk}
